\input{e.tex}

\corrPolitique{0}

\newcommand{\eps}{\varepsilon}

\begin{document}

%\maketitle

\parindent0mm

\section{Fonctions, zéros et points fixes (I)}

{\bf Motivation et exemple :} 

\medskip

On va s'intéresser au calcul approché de zéros de fonctions. Ce problème intervient dans de nombreuses applications pratiques. Prenons par exemple un fond d'investissement, sur lequel on apporte chaque année $v$ euros. Au bout de l'année numéro $n$, on retire finalement un capital de $M$ euros. On veut connaître le taux d'intérêt moyen $r$. Cela revient à trouver le zéro de la fonction $f$ suivante:
\[
f(r) = M r - {v} ( 1 + r  ) \left ( (1 + r) ^n - 1 \right ),
\]
pour laquelle on ne dispose pas de formule explicite si $n$ a une valeur trop élevée.

\bigskip

\begin{exercice}

En utilisant la relation:
\[
M = v(1+r) + \dots + v(1+r)^n,
\]
montrer qu'effectivement, le taux d'intérêt est obtenu en résolvant 
\[
f(r) = 0,
\]
avec $f$ la fonction dont l'expression est donnée ci-dessus.
Pour quelles valeurs de $n$ pouvez-vous trouver explicitement $r$ ?
\end{exercice}


\bigskip

Pour la suite, nous aurons besoin du résultat fondamental suivant:

\begin{theorem}[Théorème des valeurs intermédiaires]
Soit $f$ une fonction réelle définie et continue sur un intervalle $[a,b]$. Alors $f$ est bornée sur $[a,b]$ et pour tout réel $y$ tel que:
\[
\inf_{x \in [a,b]} f(x) \leq y \leq \sup_{x \in [a,b]} f(x),
\]
il existe (au moins) un réel $\xi \in [a,b]$ tel que 
\[
y = f(\xi).
\]
En particulier, la borne inférieure et la borne supérieure de $f$ sont atteintes et peuvent être remplacées par $\min_{x \in [a,b]}$ et $\max_{x \in [a,b]}$, respectivement.
\end{theorem}

On rappelle également une variante utile :

\begin{theorem}[Théorème de la bijection]
On se place dans le même cadre que le théorème précédent, et on suppose de plus que $f$ est strictement monotone sur $]a,b[$. Alors pour tout réel $y$ tel que:
\[
\inf_{x \in [a,b]} f(x) \leq y \leq \sup_{x \in [a,b]} f(x),
\]
il existe un et un seul réel $\xi \in [a,b]$ tel que $y = f(\xi)$.
\end{theorem}

\begin{exercice}

Soit $f$ une fonction réelle définie et continue sur $[a,b]$. On suppose que 
\[
f(a) f(b) \leq 0.
\]
Montrer qu'il existe (au moins) un réel $\xi \in [a,b]$ tel que $f(\xi)=0$.
\end{exercice}

\medskip

\begin{exercice}

On considère la fonction $f$ donnée par l'expression:
\[
f(x) = 2 - x + \ln(x).
\]

\begin{enumerate}

\item Quel est l'ensemble de définition de $f$ ?

\item Représenter graphiquement $f$.

\item Montrer que l'équation $f(x)=0$ admet exactement deux solutions réelles, qui seront notées $\alpha$ et $\beta$.

\item Montrer que
\[
\frac1{10} < \alpha < \frac12.
\]

\item Montrer que
\[
3 < \beta < 4.
\]

\item Est-ce que $3 < \beta < 3.5$ ou $3.5 < \beta < 4$ ?

\item Trouver une valeur approchée de $\beta$ à $\frac1{16}$ près.

\end{enumerate}

\end{exercice}

\medskip

\begin{exercice}

Soit la fonction $f$ donnée par l'expression:
\[
f(x) = x^3 + 2 x + 1.
\]

\begin{enumerate}

\item Quel est l'ensemble de définition de $f$ ?

\item Représenter graphiquement $f$.

\item Montrer que l'équation $f(x)=0$ admet exactement une solution réelle $\alpha$.

\item Montrer que
\[
-2 < \alpha < 0.
\]

\item Soit 
\[
h(x) = x - \frac{f(x)}{f'(x)}.
\]
Soit $x_0 = 0$ et la suite définie par récurrence :
\[
x_{n+1} = h ( x_n ), \quad n \in \eN.
\]
Calculer $x_1,x_2,x_3,\ldots$. Qu'observez-vous ?

\end{enumerate}

\end{exercice}

\medskip

\begin{exercice}

Soit $g$ une fonction réelle, définie et continue sur un intervalle $[a,b]$, telle que
\[
\forall x \in [a,b],\quad g(x) \in [a,b].
\]
Montrer qu'il existe un réel $\xi \in [a,b]$ tel que 
\[
g(\xi) = \xi.
\]
On dira que \emph{$\xi$ est un point fixe de la fonction $g$}. Le résultat démontré est d'ailleurs connu sous le nom du \emph{Théorème de Point Fixe de Brouwer}.

\medskip

Peut-on avoir plusieurs points fixes ? Si oui, donner un exemple explicite.

\end{exercice}

\newpage

\section{Fonctions, zéros et points fixes (II)}

\begin{exercice}
\label{exo:fonctionsfg}

Soit la fonction $f$ définie sur l'intervalle $[1,2]$ par l'expression
\[
f(x) = e^x - 2 x - 1, \quad x\in[1,2].
\]

\begin{enumerate}
\item Montrer qu'il existe $\xi \in [1,2]$ tel que $f(\xi)=0$.

\item Montrer que l'équation $f(x)=0$ est équivalente sur $[1,2]$ à l'équation $g(x) = x$, avec 
\[
g(x) = \ln ( 2x + 1).
\]

\item Montrer, en utilisant le Théorème de Brouwer, que $g$ admet au moins un point fixe sur $[1,2]$.

\item On choisit cette fois $g(x)=\frac12 ( e^x - 1)$. Montrer que l'équation $f(x)=0$ est encore équivalente sur $[1,2]$ à l'équation $g(x)=x$. Peut-on encore ici appliquer le Théorème de Brouwer pour montrer que $g$ admet au moins un point fixe ?

\end{enumerate}

\end{exercice}

\medskip

\begin{exercice}
\label{exo:pointFixeBanach}

On considère une fonction $g$ définie sur un intervalle $[a,b]$ et telle que pour tout $x \in [a,b]$, $g(x) \in [a,b]$.
On va supposer que $g$ est une \emph{contraction} sur $[a,b]$, autrement dit qu'il existe une constante $L$ telle que $0 < L < 1$ et que
\[
\forall x,y \in [a,b], \quad
| g(x) - g(y) | \leq L | x - y |.
\]

\begin{enumerate}
\item Montrer que $g$ est nécessairement continue.

\item Montrer que $g$ admet un \emph{unique} point fixe $\xi$ sur $[a,b]$.
\end{enumerate}

Pour $x_0 \in [a,b]$, on définit maintenant une suite $(x_k)$ par la relation de récurrence suivante:
\[
\forall k \in \eN, \quad x_{k+1} = g(x_k).
\]

\begin{enumerate}
\item[3] Montrer que la suite $(x_{k})$ tend vers $\xi$ lorsque $k$ tend vers l'infini.
\end{enumerate}

On vient de démontrer ici une version simplifiée du \emph{Théorème du Point Fixe de Banach (ou de Picard)}.
Vous verrez ultérieurement une démonstration plus générale. En particulier, on peut se passer de l'hypothèse $g([a,b]) \subset [a,b]$.

\end{exercice}

\medskip

\begin{exercice}
\label{exo:fonctiongPF}

On reprend la fonction $g$ de la question 2 de l'exercice \ref{exo:fonctionsfg}, définie sur $[1,2]$ par 
\[
g(x) = \ln(2x+1).
\]

\begin{enumerate}

\item Soit $x,y\in[1,2]$. Montrer qu'il existe $\eta \in [x,y]$ tel que:
\[
| g(x) - g(y) | = |g'(\eta)| | x - y |.
\]

\item Montrer que $g'$ est une fonction strictement décroissante sur $[1,2]$ et en déduire:
\[
\frac25 \leq g'(\eta) \leq \frac23.
\]

\item Montrer alors que $g$ est une contraction, autrement dit qu'il existe $0 < L < 1$ tel que pour $x,y\in[1,2]$,
\[
| g(x) - g(y) | \leq L | x - y |.
\]
Quelle valeur pouvez-vous donner à $L$ ?

\item Montrer que $g$ admet un unique point fixe $\xi$ sur $[1,2]$, et que la suite 
\[
x_{k+1} = \ln ( 2 x_k + 1), \quad k\in\eN,
\]
converge vers $\xi$, pour n'importe qu'elle valeur de $x_0$ dans $[1,2]$.

\item Avec le choix $x_0 = 1$, calculer $x_k$ pour $k=1,\ldots,11$. Qu'observez-vous ?

\end{enumerate}

\end{exercice}

\medskip

\begin{exercice}

On se place dans le cadre de l'exercice \ref{exo:pointFixeBanach}
et on considère toujours une fonction $g$ définie sur $[a,b]$ et telle que pour tout $x \in [a,b]$, $g(x) \in [a,b]$.
On suppose encore que $g$ est une \emph{contraction} sur $[a,b]$, autrement dit qu'il existe une constante $L$ telle que $0 < L < 1$ et que
\[
\forall x,y \in [a,b], \quad
| g(x) - g(y) | \leq L | x - y |.
\]
On note toujours $\xi$ l'unique point fixe de $g$, et pour $x_0 \in [a,b]$, on définit la suite:
\[
\forall k \in \eN, \quad x_{k+1} = g(x_k).
\]

\begin{enumerate}
\item Montrer d'abord que 
\[
| x_0 - \xi | \leq \frac1{1-L} | x_0 - x_1 |.
% | x_0 - \xi | \leq | x_0 - x_1 | + L | x_0 - \xi |.
\]

\item En déduire que pour tout $k \in \eN$, on vérifie
\[
| x_k - \xi | \leq \frac{L^k}{1-L} | x_1 - x_0 |.
\]

\item On se donne une tolérance $\varepsilon > 0$ pour l'approximation du point fixe $\xi$ par la suite $(x_k)$. On souhaite déterminer la valeur de l'itération $k$ à partir de laquelle 
\[
| x_k - \xi | \leq \varepsilon.
\]
Démontrer que l'inégalité ci-dessus est vérifiée pour tout $k$ tel que
\[
 k \geq \frac{\ln |x_1 - x_0 | - \ln (\varepsilon(1-L))}{\ln(1/L)}.
\]

\item \emph{Application:} Reprenons la fonction $g$ donnée à l'exercice \ref{exo:fonctiongPF}. Combien faut-il d'itérations pour obtenir $\xi$ avec une précision $\varepsilon = 0.5 \times 10^{-6}$ ?

\end{enumerate}

\end{exercice}

\newpage

\section{Points fixes}

On va commencer par donner une définition qui permet de caractériser différents comportements des points fixes d'une fonction:

\begin{definition}
 Soit $g$ une fonction réelle définie sur un intervalle $I$. Supposons que $g$ admette (au moins) un point fixe, noté $\xi\in I$, et que $g$ soit dérivable en $\xi$. Alors :
 \begin{itemize}
  \item si $|g'(\xi)| < 1$, on dira que $\xi$ est un point fixe \emph{stable.}
  \item si $|g'(\xi)| > 1$, on dira que $\xi$ est un point fixe \emph{instable.}
 \end{itemize}

\end{definition}

On commence par donner le résultat suivant, qui peut être vu comme une version locale du Théorème de point fixe de Banach.

\begin{theorem}[Convergence au voisinage d'un point fixe stable]
\label{thm:pointfixelocal}
Soit $g$ une fonction réelle définie et continue sur un intervalle $I$. On suppose qu'il existe 
%\[
%\forall x \in [a,b],\, g(x) \in [a,b].
%\]
$\xi (= g(\xi)) \in I$ qui est un point fixe de $g$. On supposera que $g$ admet une dérivée continue dans un voisinage de $\xi$, et que de plus
\[
| g'(\xi) | < 1.
\]
Alors la suite $(x_k)_{k\in\eN}$ donnée par la relation de récurrence 
$$x_{k+1} = g(x_k),\quad k\geq0$$ 
converge vers $\xi$ lorsque $k\rightarrow +\infty$ si la première valeur $x_0$ est choisie suffisamment proche de $\xi$.
\end{theorem}

%Dans le cadre de ce théorème, on parlera de point fixe \emph{stable}.

\bigskip

\begin{exercice}

Démontrer le Théorème \ref{thm:pointfixelocal}.

\end{exercice}

\medskip

\begin{exercice}

On se place toujours sous les hypothèses du Théorème \ref{thm:pointfixelocal}. On veut connaître la \emph{vitesse de convergence} de la suite $(x_k)_{k\in\eN}$ qui approche le point fixe $\xi$. Montrer que
\[
\lim_{k\rightarrow +\infty} \frac{|x_{k+1} - \xi|}{| x_k - \xi |} = |g'(\xi)|.
\]

\end{exercice}

\medskip

\begin{exercice}

Soit un réel $\alpha \in ]0;+\infty[$, on définit sur l'intervalle $[0,1]$ une fonction $g$ dépendante de $\alpha$ à l'aide de l'expression suivante:
\[
g(x) = \left\{
\begin{array}{ll}
2^{-\{ 1 + (\log_2(1/x))^{\frac1\alpha} \}^\alpha}&\quad \mbox{si } x\in]0;1],\\
0&\quad \mbox{sinon}.
\end{array}
\right.
\]

\begin{enumerate}

\item Calculer la limite de $g$ en $0+$. Que peut-on en déduire sur la continuité de $g$ ?

\item Etudier les variations de $g$ sur $[0,1]$. Représenter $g$ pour $\alpha = \frac12,1,2$.

\item La fonction $g$ admet-elle un (ou des) point(s) fixe(s) ? Le(s)quel(s) ?

\item Soit la suite $(x_k)_{k\in\eN}$ donnée par $x_0=1$ et $x_{k+1} = g(x_k)$. Montrer que pour tout $k\geq0$, nous avons
\[
x_{k} = 2^{-k^\alpha}.
\]
Qu'en déduit-on sur la convergence de la suite ?

\item Calculer la limite suivante, qui va donner la vitesse de convergence de la suite, en fonction de $\alpha$:
\[
\lim_{k \rightarrow +\infty} \frac{|x_{k+1}|}{|x_k|}.
\]
Comment interprète-t'on ce résultat ?

\end{enumerate}

\end{exercice}

\medskip

\begin{exercice}

Soit $g$ une fonction réelle définie et continue sur un intervalle $I$. On suppose que $g$ admet un point fixe sur $I$,
%\[
%\forall x \in [a,b],\, g(x) \in [a,b].
%\]
qu'on notera $\xi (= g(\xi)) \in I$. %un point fixe de $g$, et 
On supposera que $g$ admet une dérivée continue dans un voisinage de $\xi$, et que de plus
\[
| g'(\xi) | > 1.
\]
Autrement dit $\xi$ est un point fixe instable.
On se donne toujours la suite $(x_k)_{k\in\eN}$ avec la relation de récurrence 
$$x_{k+1} = g(x_k),\quad k\geq0.$$

Montrer que, pour tout $x_0$ dans $I$, % si $x_0$ est suffisamment proche de $\xi$, %avec $x_0 \neq \xi$, 
alors, lorsque $k\rightarrow +\infty$:
\begin{enumerate}
 \item soit la suite $x_k$ atteint exactement $\xi$ après un nombre fini d'itérations $k_0 \geq 0$ ($x_{k} = \xi$ pour tout $k \geq k_0$).
 \item soit la suite $(x_k)_{k\in\eN}$ ne converge pas vers $\xi$.
\end{enumerate}

%\smallskip

%On dira que le point fixe $\xi$ est \emph{instable}.

\end{exercice}

\medskip

\begin{exercice}

Soit la fonction $g$ définie sur $\eR$ à l'aide de l'expression suivante:
\[
g(x) = \frac12 ( x^2 + c ),
\]
où $c\in\eR$ est une constante donnée.

\begin{enumerate}

\item Montrer qu'en fonction de $c$, $g$ admet $0$, $1$ ou $2$ points fixes. Donner leur expression (en fonction de $c$).

\item On se place maintenant dans le cadre où $g$ admet deux points fixes $\xi_1$ et $\xi_2$, avec la convention $\xi_1 < \xi_2$. Montrer que 
\[
\xi_1 < 1 < \xi_2.
\]

\item Montrer que $\xi_2$ est instable, autrement dit que $|g'(\xi_2)| > 1$.

\item Est-ce que $\xi_1$ est stable ($|g'(\xi_1)| < 1$) ? Pour quelles plages de valeurs de $c$ ?

\item On va supposer ici $0 < c < 1$ et s'intéresser à la suite $(x_k)$ donnée par $x_{k+1} = g(x_k)$. Montrer que 
\[
x_{k+1} - \xi_1 = \frac12 ( x_k + \xi_1 ) ( x_k - \xi_1 ), \quad k\geq 0.
\]
En déduire que
\[
\lim_{k \rightarrow +\infty} x_k = \xi_1
\]
pour $0 \leq x_0 < \xi_2$.

\item Que peut-on dire sur la vitesse de convergence 
\[
\lim_{k\rightarrow +\infty} \frac{|x_{k+1} - \xi_1|}{| x_k - \xi_1 |}
\]
lorsque $c=0$ et lorsque $c$ est proche de $1$ ?

\end{enumerate}

\end{exercice}

\newpage

\section{Méthodes de relaxation et de Newton-Raphson}

On supposera ici que $f$ est une fonction réelle définie et continue sur un voisinage d'un réel noté $\xi$. On appelle \emph{méthode de relaxation} la méthode qui permet d'approcher un zéro de $f$ à l'aide de la suite $(x_k)$ définie par la relation de récurrence
\begin{equation} \label{eq:relaxation}
x_{k+1} = x_k - \lambda f(x_k), \quad k\geq0,
\end{equation}
où le réel $\lambda\neq0$ est un paramètre fixé, choisi judicieusement, et où la valeur initiale $x_0$ est donnée dans un voisinage de $\xi$.

\medskip

Remarquons que cela correspond à faire des itérations de point fixe pour une fonction $g$ définie par
\[
g(x) = x - \lambda f(x).
\]

\medskip

\begin{exercice}

On va établir ici un résultat de convergence pour la méthode de relaxation, moyennant les bonnes conditions sur $f$ et sur $\lambda$. 
On va pour cela supposer que $f(\xi)=0$, que $f'$ est définie et continue sur un voisinage de $\xi$, et que $f'(\xi)\neq0$.

Montrer alors qu'il existe $\delta > 0$, tel que pour tout $\lambda$ suffisamment petit (et $\lambda \neq 0$), la suite $(x_k)$ définie par \eqref{eq:relaxation} converge vers $\xi$ pour tout $x_0$ dans l'intervalle $[\xi - \delta ; \xi + \delta ]$.

\end{exercice}

\medskip

\begin{exercice}

On modifie ici la méthode de relaxation \eqref{eq:relaxation} en s'autorisant un choix différent de $\lambda$ à chaque itération:
\[
x_{k+1} = x_k - \lambda(x_k) f(x_k), \quad k\geq0.
\]
On supposera que $\lambda$ est une fonction continue par rapport à $x$.

\begin{enumerate}

\item Montrer que si la suite $(x_k)$ converge vers un réel $\xi$, avec $\lambda(\xi)\neq0$, alors $\xi$ est solution de $f(\xi)=0$.

\item Montrer que, si $f$ est dérivable en $\xi$ et si $(x_k)$ converge vers $\xi$, la vitesse de convergence asymptotique de la suite est donnée par 
\[
\lim_{k\rightarrow +\infty} \frac{|x_{k+1} - \xi|}{| x_k - \xi|}
= | 1 - \lambda(\xi) f'(\xi)|.
\]

\item Quel choix paraît alors judicieux pour $\lambda(x)$ ?

\end{enumerate}

\end{exercice}

\medskip

On supposera ici que $f$ est dérivable dans un voisinage de $\xi$, avec de plus sa dérivée $f'$ qui ne s'annule pas dans ce voisinage.
On définit maintenant la \emph{méthode de Newton-Raphson} avec la relation de récurrence
\begin{equation} \label{eq:newton}
x_{k+1} = x_k - \frac{f(x_k)}{f'(x_k)}, \quad k\geq0,
\end{equation}
où la valeur initiale $x_0$ est donnée dans un voisinage de $\xi$.

\medskip

Remarquons que cela correspond aussi à faire des itérations de point fixe pour une fonction $g$ définie par
\[
g(x) = x - \frac{f(x)}{f'(x)}.
\]

\medskip

\begin{exercice}

Nous allons établir ici un résultat de convergence pour la méthode de Newton-Raphson. Nous supposerons pour cela que $f$ est deux fois dérivable, avec une dérivée seconde $f''$ continue sur un intervalle $I_\delta = [ \xi - \delta ; \xi + \delta ]$, $\delta > 0$. On supposera de plus $f(\xi) = 0$, $f''(\xi)\neq0$, et qu'il existe $A >0$ tel que
\begin{equation} \label{eq:A}
\left | \frac{f''(x)}{f'(y)} \right | \leq A, \quad \forall x,y \in I_\delta.
\end{equation}
On notera par ailleurs $h = \min \{ \delta , \frac1{A} \}$.

\begin{enumerate}

\item Supposons qu'à l'itération $k\geq0$, le réel $x_k$ soit tel que $|x_k - \xi| \leq h$. Montrer qu'il existe $\eta_k \in I_\delta$ tel que 
\[
0 = f(x_k) + (\xi - x_k) f'(x_k) + \frac{(\xi-x_k)^2}{2}f''(\eta_k).
\]

\item En déduire que
\[
\xi - x_{k+1} = - \frac{(\xi-x_k)^2 f''(\eta_k)}{2 f'(x_k)},
\]
puis que
\[
| \xi - x_{k+1} | \leq \frac12 | \xi - x_k |.
\]

\item Montrer alors que, si $x_0$ est tel que $|x_0 - \xi|\leq h$, nous avons pour tout $k\geq0$:
\[
| \xi - x_k | \leq h 2^{-k}.
\]
Que peut-on en déduire ?

\item Montrer maintenant que
\[
\lim_{k\rightarrow +\infty} \frac{|x_{k+1} - \xi|}{| x_k - \xi|^2}
= \mu.
\]
Donner l'expression du réel $\mu$ et montrer que
\[
0 < \mu \leq \frac{A}{2}.
\]

\medskip

\end{enumerate}

Nous avons finalement établi que la méthode de Newton converge avec un taux de convergence quadratique.

\end{exercice}

\medskip

\begin{exercice}

%On se place dans le même cadre que l'exercice précédent, avec les mêmes hypothèses sur $f$. 
Soit $f$ une fonction réelle deux fois dérivable sur un intervalle $I$, avec une dérivée seconde $f''$ continue sur $I$. On supposera de plus qu'il existe $\xi \in I$ avec $f(\xi) = 0$, $f''(\xi)\neq0$,
et qu'il existe un réel $X > \xi$ tel que sur l'intervalle $J=[\xi,X] \subset I$, les dérivées première $f'$ et seconde $f''$ de $f$ soient strictement positives. Montrer alors que la suite $(x_k)$ définie en \eqref{eq:newton} converge vers $\xi$, avec un taux quadratique, pour toute valeur initiale $x_0$ dans $J$.

\end{exercice}

\medskip

\begin{exercice}

On reprend ici l'exercice 1.3, avec la fonction $f$ donnée par l'expression:
\[
f(x) = x^3 + 2 x + 1.
\]
On rappelle que cette fonction admet un unique zéro $\alpha$ avec
\[
-2 < \alpha < 0.
\]

\begin{enumerate}

\item Calculer une valeur approchée de $\alpha$ en utilisant la méthode de relaxation \eqref{eq:relaxation}. Comment choisissez-vous le paramètre de relaxation $\lambda$ et combien d'itérations faut-il alors pour obtenir une précision acceptable ?

\item Comparer les résultats obtenus à la question 1 avec ceux obtenus en utilisant la méthode de Newton-Raphson \eqref{eq:newton}. Que constatez-vous ?

\item Comment vérifie-t'on en pratique qu'il y a une convergence quadratique ? 

\end{enumerate}

\end{exercice}

\medskip

\begin{exercice}

On revient ici à l'application en économétrie. On rappelle que le taux d'intérêt $r$ pour le fond d'investissement considéré est donné par l'équation:
\[
f(r)=0, \quad f(r) = M r - v {(1+r)} \left ( (1+r)^n - 1 \right).
\]

On choisira ici : un capital initial $v=1000$ euros, le capital retiré $M=6000$ euros, au bout de $n=5$ années. Trouver une valeur approchée de $r$ à l'aide de la méthode de votre choix, et commenter sur la qualité du résultat obtenu.

\end{exercice}

\newpage

\section{Interpolation de Lagrange}

\emph{Motivation :} On mesure la température à l'extérieur du bâtiment "Mirande" pendant toute une journée:

\medskip

\begin{center}
\begin{tabular}{c|c|c|c|c|c}
$t$ & 8h & 9h & 11h & 13h & 15h \\ \hline
$T$ ($^o$C) & 12.1 & 13.6 & 15.9 & 18.5 & 16.1 
\end{tabular}
\end{center}

On veut pouvoir trouver une expression explicite de la fonction température $T(t)$, qui permettra de donner par exemple la température qu'il faisait à 10h.

\medskip

\emph{Problème général :} Soit $n\in\eN$ et $\mathcal{P}_n$ l'ensemble des polynômes sur $\eR$  de degré maximal $n$. Soit également $n+1$ réels $x_0,x_1,\ldots,x_n$ \emph{distincts} et $n+1$ réels $y_0,y_1,\ldots,y_n$. On veut trouver $p_n$ dans $\mathcal{P}_n$ tel que
\begin{equation} \label{eq:interp}
p_n(x_i)=y_i,\quad i=0,\ldots,n.
\end{equation}

\medskip

\begin{exercice} 

On supposera ici que $n\geq1$. 

\begin{enumerate}

\item Montrer qu'il existe $n+1$ polyn\^omes $L_k \in \mathcal{P}_n$, $k=0,\ldots,n$ tels que
\[
L_k(x_i) = \left \{ \begin{array}{ll}
1 & \quad i=k,\\
0 & \quad i\neq k,
\end{array} \right.
\]

L'ensemble de ces polynômes s'appelle la \emph{base de Lagrange}.

\item Comment trouve-t'on à partir de la question 1 un polynôme qui vérifie \eqref{eq:interp} ?

\item Soit maintenant deux polynômes $p_n$ et $q_n$ qui vérifient \eqref{eq:interp}. Montrer que $p_n=q_n$.

\end{enumerate}

Lorsque $n=0$, montrer qu'il existe un unique polynôme $p_0$ qui vérifie \eqref{eq:interp}.

\end{exercice}

\medskip

L'exercice 5.1 a permis d'établir le résultat suivant:

\begin{theorem}[Théorème d'interpolation de Lagrange]
Soit $n\in\eN$. Soit $n+1$ réels $x_0,x_1,\ldots,x_n$ \emph{distincts} et $n+1$ réels $y_0,y_1,\ldots,y_n$. Il existe un unique polynôme $p_n \in \mathcal{P}_n$ tel que
\begin{equation*} %\label{eq:interp}
p_n(x_i)=y_i,\quad i=0,\ldots,n.
\end{equation*}
Ce polynôme est appelé \emph{polynôme d'interpolation de Lagrange}. Les valeurs $x_i$ sont appelées des \emph{points d'inter\-polation}.

\end{theorem}

\medskip

\begin{exercice}

Soit $f$ la fonction définie sur $[-1,1]$ à l'aide de l'expression suivante:
\[
f(x) = e^x.
\]
Trouver le polynôme d'interpolation de Lagrange associé à $f$ sur $[-1,1]$, en prenant les points d'interpolation suivants: $x_0=-1$, $x_1=0$, $x_2=1$ (et donc $y_i = f(x_i)$, $i=0,1,2$).

\end{exercice}

\medskip

\begin{exercice}

Le but de cet exercice est d'estimer \emph{l'erreur d'interpolation}, autrement dit la distance entre une fonction $f$ et un polynôme d'interpolation de Lagrange $p_n$ défini à partir de valeurs de $f$. On suppose pour cela que $f$ est une fonction réelle définie et continue sur un intervalle $[a,b]$, qui admet des dérivées continues sur $[a,b]$ jusqu'à l'ordre $n+1$. On prendra toujours $n+1$ points d'interpolation distincts $x_0,x_1,\ldots,x_n$ et on notera $p_n \in \mathcal{P}_n$ le polynôme d'interpolation associé ($p_n(x_i) = f(x_i)$, $i=0,\ldots,n$).

On définit
\[
\pi_{n+1} (x) = (x - x_0) (x - x_1) \ldots (x - x_n).
\]

On prend $x\in[a,b]$, $x\neq x_i$, $i=0,\ldots,n$. On introduit alors la fonction auxilliaire
\[
\varphi(t) = f(t) - p_n(t) - \frac{f(x) - p_n(x)}{\pi_{n+1}(x)} \pi_{n+1}(t).
\]

\begin{enumerate}

\item Justifier que $\varphi$ est bien définie sur $[a,b]$.

\item Montrer que $\varphi'$ s'annule en $n+1$ points sur l'intervalle $]a;b[$.

\item Dans le cas $n=0$, montrer alors qu'il existe $\xi \in ]a;b[$ tel que 
\[
f(x) - p_0(x) = f'(\xi) \pi_1(x).
\]

\item On suppose maintenant $n\geq1$.

\begin{enumerate}

\item Montrer qu'il existe $\xi \in ]a,b[$ tel que $\varphi^{(n+1)}(\xi) = 0$.

\item En déduire

\[
0 = f^{(n+1)} (\xi) - \frac{f(x) - p_n(x)}{\pi_{n+1}(x)} (n+1)!
\]

\end{enumerate}

\item Déduire de ce qui précède que, pour $n\geq0$, il existe $\xi \in ]a;b[$ tel que 
\[
f(x) - p_n(x) = \frac{f^{(n+1)}(\xi)}{(n+1)!} \pi_{n+1}(x).
\]

\item Etablir finalement l'inégalité suivante:
\[
| f(x) - p_n(x) | \leq \frac{M_{n+1}}{(n+1)!} | \pi_{n+1}(x) |,
\]
où on précisera la valeur de la constante $M_{n+1} \geq 0$

\end{enumerate}

\end{exercice}

\newpage

\section{Interpolation de Hermite}

\begin{theorem}[Théorème d'interpolation de Hermite]
Soit $n\in\eN$. Soit $n+1$ réels $x_0,x_1,\ldots,x_n$ \emph{distincts}. Soit également deux ensembles de $n+1$ réels: d'une part $y_0,y_1,\ldots,y_n$ et d'autre part $z_0,z_1,\ldots,z_n$. Il existe un unique polynôme $p_{2n+1} \in \mathcal{P}_{2n+1}$ tel que
\begin{equation} \label{eq:interpH}
p_{2n+1}(x_i)=y_i,\quad 
p_{2n+1}'(x_i)=z_i,\quad i=0,\ldots,n.
\end{equation}
Ce polynôme est appelé \emph{polynôme d'interpolation de Hermite}.
\end{theorem}

\begin{exercice}

Le but de cet exercice est de démontrer le Théorème d'interpolation de Hermite.
On considère pour cela les polynômes auxilliaires $H_k$ et $K_k$, $k=0,\ldots,n$, définis comme suit:
\begin{eqnarray*}
H_k(x) & = & [L_k(x)]^2 ( 1 - 2 L_k'(x_k)(x-x_k)),\\
K_k(x) & = & [L_k(x)]^2 ( x - x_k ),
\end{eqnarray*}
où les $L_k$ sont les polynômes de la base de Lagrange définis lors du TD précédent.
On se place d'abord dans le cas $n\geq 1$.

\begin{enumerate}

\item Quel est le degré des polynômes $H_k$ et $K_k$ ?

\item Montrer que pour tous les indices $i,k=0,\ldots,n$:
\[
H_k(x_i) = \left \{ \begin{array}{ll} 1,& i=k, \\ 0,&i\neq k, \end{array} \right.
\quad
H_k'(x_i) = 0.
\]
\item Montrer aussi que pour tous les indices $i,k=0,\ldots,n$:
\[
K_k(x_i) = 0,
\quad
K_k'(x_i) = \left \{ \begin{array}{ll} 1,& i=k, \\ 0,&i\neq k. \end{array} \right.
\]

\item Montrer alors qu'il existe un polynôme $p_{2n+1} \in \mathcal{P}_{2n+1}$ qui vérifie les conditions \eqref{eq:interpH}.

\item Montrer que s'il existe deux polynômes $p_{2n+1} \in \mathcal{P}_{2n+1}$ et $q_{2n+1} \in \mathcal{P}_{2n+1}$, alors $p_{2n+1} = q_{2n+1}$.

\item Que peut-on dire dans le cas $n=0$ ?

\end{enumerate}

\end{exercice}

\medskip

\begin{exercice}

Construire le polynôme $p_3$ de degré 3 tel que
\[
p_3(0) = 0, \quad p_3(1)=1, \quad p_3'(0)=1, \quad p_3'(1)=0.
\]
Le représenter graphiquement.

\end{exercice}

\medskip

\begin{exercice}

Soit $n\geq0$ et soit $f$ une fonction réelle, définie, continue et $2n+2$ fois dérivable sur un intervalle $[a,b]$, avec de plus $f^{(2n+2)}$ continue sur $[a,b]$.
%
Soit $n+1$ réels $x_0,x_1,\ldots,x_n$ \emph{distincts}, et $p_{2n+1} \in \mathcal{P}_{2n+1}$ le \emph{polynôme d'interpolation de Hermite pour $f$}, c'est à dire le polynôme tel que
\begin{equation*} %\label{eq:interpH}
p_{2n+1}(x_i)=f(x_i),\quad 
p_{2n+1}'(x_i)=f'(x_i),\quad i=0,\ldots,n.
\end{equation*}

\begin{enumerate}

\item Montrer que pour tout $x\in[a,b]$, il existe $\xi \in ]a,b[$ tel que
\[
f(x) - p_{2n+1} (x) = \frac{f^{(2n+2)}(\xi)}{(2n+2)!} [\pi_{n+1}(x)]^2,
\]
où le polynôme $\pi_{n+1}$ a été défini lors du TD précédent.

\item En déduire l'inégalité suivante:
\[
| f(x) - p_{2n+1} (x) | \leq \frac{M_{2n+2}}{(2n+2)!} [\pi_{n+1}(x)]^2,
\]
où on précisera la valeur de la constante $M_{2n+2} \geq 0$

\end{enumerate}

\end{exercice}

\newpage

\section{Intégration et quadratures (I)}

\begin{exercice}

Calculer les intégrales suivantes:

\begin{enumerate}

\item D'abord $$ \int_0^1 e^x dx \qquad \mbox{ et } \qquad \int_0^\pi \cos(x) dx.$$

\item Ensuite $$ \int_0^1 e^{x^2} dx \qquad \mbox{ et } \qquad \int_0^\pi \cos(x^2) dx.$$

\item Et finalement
\[
\int_1^{2000} \exp(\sin(\cos(\sinh(\cosh(\arctan(\log(x))))))) dx.
\]

\end{enumerate}

\end{exercice}

\medskip

\begin{exercice}

Dans cet exercice, nous allons établir la \emph{formule des trap\`ezes}, qui permet de calculer une intégrale de façon approchée. On l'appelle aussi \emph{formule de Newton-Cotes d'ordre 1}. On considère pour cela une fonction réelle $f$ définie et continue sur un intervalle $[a,b]$. Soit $p_1$ le polynôme de Lagrange de degré 1 qui interpole $f$ aux points $a$ et $b$. On rappelle qu'on peut écrire $p_1$ comme suit:
\[
p_1(x) = L_0(x) f(a) + L_1(x) f(b),
\]
avec $L_0$ et $L_1$ qui sont les polynômes de degré 1 de la base de Lagrange.

\begin{enumerate}

\item Montrer que:
\[
\int_a^b p_1(x) dx = \frac{b-a}2 ( f(a) + f(b) ).
\]

L'idée est ensuite de remplacer $f$ par $p_1$ dans le calcul d'intégrale, ce qui donne la formule approchée 
\begin{equation} \label{eq:trapezes}
\int_a^b f(x) dx \simeq \frac{b-a}2 ( f(a) + f(b) ).
\end{equation}

\item Appliquer la formule \eqref{eq:trapezes} pour calculer de façon approchée les intégrales suivantes:

$$ \int_0^1 e^x dx \qquad \mbox{ et } \qquad \int_0^\pi \cos(x) dx.$$

$$ \int_0^1 e^{x^2} dx \qquad \mbox{ et } \qquad \int_0^\pi \cos(x^2) dx.$$

Comparer vos résultats à ceux de l'exercice 1.

\end{enumerate}

\end{exercice}

\medskip

\begin{exercice}

De façon similaire à l'exercice précédent, nous allons établir la \emph{formule de Simpson}, toujours dans le but de calculer une intégrale de façon approchée (on parle encore de \emph{formule de Newton-Cotes d'ordre 2}). Soit une fonction réelle $f$ définie et continue sur un intervalle $[a,b]$. Soit $p_2$ le polynôme de Lagrange de degré 2 qui interpole $f$ aux points $a$, $m=\frac{a+b}2$ et $b$.

\begin{enumerate}

\item Montrer que:
\[
\int_a^b p_2(x) dx = \frac{b-a}6 ( f(a) + 4f(m) + f(b) ).
\]

On obtient comme cela la formule approchée 
\begin{equation} \label{eq:Simpson}
\int_a^b f(x) dx \simeq \frac{b-a}6 ( f(a) + 4f(m) + f(b) ).
\end{equation}

\item Appliquer la formule \eqref{eq:Simpson} pour calculer de façon approchée les intégrales suivantes:

$$ \int_0^1 e^x dx \qquad \mbox{ et } \qquad \int_0^\pi \cos(x) dx.$$

$$ \int_0^1 e^{x^2} dx \qquad \mbox{ et } \qquad \int_0^\pi \cos(x^2) dx.$$

Comparer vos résultats à ceux des exercices 1 et 2.

\end{enumerate}

\end{exercice}

\medskip

\begin{exercice}

Soit $f$ une fonction réelle définie et continue sur un intervalle $[a,b]$. Soit $h=(b-a)/n$, avec $n\geq1$, et soient les points d'interpolation définis comme suit
\[
x_i = a + i h, \qquad i=0,\ldots,n.
\]
Soit $p_n \in \mathcal{P}_n$ le polynôme d'interpolation de $f$ aux points $x_i$.

\begin{enumerate}

\item Montrer que
\[
 \int_a^b p_n(x) dx = \sum_{k=0}^n w_k f(x_k),
\]
et donner une expression des $w_k$.

\item Soit \emph{l'erreur de quadrature} la différence entre la valeur exacte de l'intégrale et sa valeur approchée en utilisant la formule de la question 1:
\[
E_n (f) = \int_a^b f(x) dx - \sum_{k=0}^n w_k f(x_k).
\]
Montrer que l'on peut majorer cette erreur de la façon qui suit :
\[
| E_n (f) | \leq \frac{M_{n+1}}{(n+1)!} \int_a^b | \pi_{n+1} (x)| dx,
\]
avec
$M_{n+1} = \max_{\eta \in [a,b]} | f^{(n+1)} (\eta) |$
et
$\pi_{n+1}(x) = ( x - x_0 ) \ldots (x - x_n)$.

\item En déduire la majoration de l'erreur suivante pour la \emph{méthode des trapèzes}:
\[
| E_1 (f) | \leq \frac{(b-a)^3}{12} M_2.
\]

\item En déduire également la majoration de l'erreur suivante pour la \emph{méthode de Simpson}:
\[
| E_2 (f) | \leq \frac{(b-a)^4}{196} M_3.
\]


\end{enumerate}

\end{exercice}

\newpage

\section{Intégration et quadratures (II)}

On cherche à calculer l'intégrale d'une fonction $f$ sur un intervalle $[a,b]$ quelconque. On va pour cela construire des \emph{formules de quadrature} composées en subdivisant l'intervalle $[a,b]$ en sous-intervalles de petite taille, et en appliquant les formules vues au TD 1 sur chaque sous-intervalle.

Dans la suite de ce TD, on notera $f: [a,b] \rightarrow \eR$ une fonction continue, dont on souhaite approcher l'intégrale. Une subdivision $(x_i)$ de l'intervalle $[a,b]$ sera notée comme suit :
\[
x_i = a + i h, \quad i=0,\ldots,m,
\]
où $m \in \eN$, $m\geq1$, est le nombre de points de la subdivision, et $h = (b-a)/m$ est le \emph{pas} de la subdivision.

\bigskip

\begin{exercice}

En appliquant la formule des trapèzes sur chaque intervalle $[x_i,x_{i+1}]$ de la subdivision $(x_i)$, montrer qu'on peut approcher l'intégrale de $f$ sur $[a,b]$ par la formule suivante :
\begin{equation}\label{eq:trapezes_comp}
\int_a^b f(x) dx \simeq
h \left (
\frac12 f(x_0) + f(x_1) + \ldots + f(x_{m-1}) + \frac12 f(x_m)
\right ).
\end{equation}

\end{exercice}

\medskip

\begin{exercice}

Appliquer la formule \eqref{eq:trapezes_comp} pour calculer de façon approchée les intégrales suivantes:

$$ \int_0^1 e^x dx \qquad \mbox{ et } \qquad \int_0^\pi \cos(x) dx,$$

$$ \int_0^1 e^{x^2} dx \qquad \mbox{ et } \qquad \int_0^\pi \cos(x^2) dx.$$

On choisira des subdivisions de $[0,1]$ en deux intervalles de même taille, puis en quatre intervalles de même taille.

\medskip

Comparer vos résultats à ceux du TD précédent.

\end{exercice}

\medskip

\begin{exercice}

On définit :
\[
\mathcal{E}_1 =
\int_a^b f(x) dx -
h \left (
\frac12 f(x_0) + f(x_1) + \ldots + f(x_{m-1}) + \frac12 f(x_m)
\right )
\]
l'erreur de quadrature associée à la formule composite des trapèzes. Montrer qu'on peut majorer cette erreur comme suit :
\[
| \mathcal{E}_1 | \leq \frac{(b-a)^3}{12 m^2} \max_{\zeta\in[a;b]} |f''(\zeta)|.
\]

Commenter le résultat obtenu.

\end{exercice}

\medskip

\begin{exercice}

Pour chaque indice $i=0,\ldots,m-1$, on notera 
\[
x_{i+\frac12} = \frac12 ( x_{i} + x_{i+1} ),
\]
le point au milieu du segment $[x_i;x_{i+1}]$.
En appliquant la formule de Simpson sur chaque intervalle $[x_i,x_{i+1}]$ de la subdivision $(x_i)$, montrer qu'on peut approcher l'intégrale de $f$ sur $[a,b]$ par la formule suivante :
\begin{equation}\label{eq:Simpson_comp}
\int_a^b f(x) dx \simeq
\frac{h}{6} \left (
f(x_0) + 4f(x_{\frac12}) + 2 f(x_1) + 4f(x_{\frac32}) +
\ldots + 2f(x_{m-1}) + 4 f(x_{m-\frac12}) + f(x_m)
\right ).
\end{equation}

\end{exercice}

\medskip

\begin{exercice}

Appliquer la formule \eqref{eq:Simpson_comp} pour calculer de façon approchée les intégrales suivantes:

$$ \int_0^1 e^x dx \qquad \mbox{ et } \qquad \int_0^\pi \cos(x) dx,$$

$$ \int_0^1 e^{x^2} dx \qquad \mbox{ et } \qquad \int_0^\pi \cos(x^2) dx.$$

On choisira des subdivisions de $[0,1]$ en deux intervalles de même taille, puis en quatre intervalles de même taille.

\medskip

Comparer vos résultats à ceux du TD précédent et à ceux de l'exercice 8.2.

\end{exercice}

\medskip

\begin{exercice}

On définit :
\[
\mathcal{E}_2 =
\int_a^b f(x) dx -
\frac{h}{6} \left (
f(x_0) + 4f(x_{\frac12}) + 2 f(x_1) + 4f(x_{\frac32}) +
\ldots + 2f(x_{m-1}) + 4 f(x_{m-\frac12}) + f(x_m)
\right )
\]
l'erreur de quadrature associée à la formule composite de Simpson. Montrer qu'on peut majorer cette erreur comme suit :
\[
| \mathcal{E}_2 | \leq \frac{(b-a)^4}{196 m^3} \max_{\zeta\in[a;b]} |f^{(3)}(\zeta)|.
\]

Commenter le résultat obtenu et comparer avec celui de l'exercice 8.3.

\end{exercice}

\newpage

\section{Fonctions à plusieurs variables (I)}

Le but de ce TD est d'étudier des fonctions à plusieurs variables, autrement dit de la forme :
\[
f : \left \{ 
\begin{array}{lcl}
\eR^N & \rightarrow & \eR \\
(x_1,x_2,\ldots,x_N) & \mapsto & f(x_1,x_2,\ldots,x_N)
\end{array}
\right.
\]
avec $N \geq 2$. On examinera d'abord leur continuité, puis ensuite on s'intéressera aux extrema de ces fonctions.
On rappelle que le \emph{gradient} de $f$ est le vecteur de composantes :
\[
\nabla f = \left [ \frac{\partial f}{\partial x_1} , 
\frac{\partial f}{\partial x_2} , \ldots , \frac{\partial f}{\partial x_N} \right ]^\mathrm{T}
\]
et on rappelle également que la \emph{matrice hessienne} de $f$ s'écrit :
\[
\mbox{Hess } f = \left ( \frac{\partial^2 f}{\partial x_i \partial x_j} \right )_{1\leq i,j \leq N}.
\]

\begin{exercice}

Etudier la continuité des fonctions à deux variables suivantes :

\begin{align*}
& f_1(x,y) = |x| \\ 
& f_2(x,y) = \left\{ \begin{array}{ll}
\displaystyle \frac{\sin(x^2 + y^2)}{x^2 + y^2} & \mbox{ si } (x,y)\neq(0,0)\\\\
1 & \mbox{ sinon}
\end{array} \right.\\
& f_3(x,y) = \left\{ \begin{array}{ll}
\displaystyle \frac{xy}{\sqrt{x^2 + y^2}} & \mbox{ si } (x,y)\neq(0,0)\\\\
0 & \mbox{ sinon}
\end{array} \right.\\
& f_4(x,y) = \left\{ \begin{array}{ll}
\displaystyle \frac{x^2}{\sqrt{x^2 + y^2}} & \mbox{ si } (x,y)\neq(0,0)\\\\
0 & \mbox{ sinon}
\end{array} \right.\\
& f_5(x,y) = \left\{ \begin{array}{ll}
\displaystyle \frac{x^2}{{x^2 + y^2}} & \mbox{ si } (x,y)\neq(0,0)\\\\
0 & \mbox{ sinon}
\end{array} \right.\\
& f_6(x,y) =  \left\{ \begin{array}{ll}
\displaystyle \frac{x^2 y}{x^4 + y^2} & \mbox{ si } (x,y)\neq(0,0)\\\\
0 & \mbox{ sinon}
\end{array} \right.\\
& f_7(x,y) = \left\{ \begin{array}{ll}
\displaystyle xy \frac{x^2 - y^2}{x^2 + y^2} & \mbox{ si } (x,y)\neq(0,0)\\\\
0 & \mbox{ sinon}
\end{array} \right.
\end{align*}

\end{exercice}

%\medskip
\newpage

\begin{exercice}

D\'eterminer les points critiques des fonctions \`a deux variables suivantes. Dire si chaque point est un minimum ou un maximum en justifiant.

\begin{align*}
&f_1(x,y)=x^2+y^2 \\\\
&f_2(x,y)=x^2-y^2 \\\\
&f_3(x,y)=x^2+3y^2+2xy-2x-10y+9 \\\\
&f_4(x,y)=x^3+x^2y-4y^2+1\\\\
&f_5(x,y)=e^{-(x^2+y^2)}\\\\
&f_6(x,y)=xye^{-(x^2+y^2)}
\end{align*}

\end{exercice}

\medskip

\begin{exercice}

D\'eterminer les points critiques des fonctions \`a trois variables suivantes. Dire si chaque point est un minimum ou un maximum en justifiant.

\begin{align*}
&f_1(x,y,z)=x^2+y^2+z^2\\\\
&f_2(x,y,z)=x^2-y^2+z^2\\\\
&f_3(x,y,z)=e^{-(x^2+y^2+z^2)}\\\\
&f_4(x,y,z)=(2x+y)^2 + x^3 + y^3 + z^3
\end{align*}

\end{exercice}

\medskip

\begin{exercice}

Soit $f : \eR^N \rightarrow \eR$ une fonction continue et \emph{coercive}, c'est à dire :
\[
\lim_{\| x \|_2 \rightarrow + \infty} f(x) = + \infty,
\]
avec $\| x \|_2^2 = x_1^2 + \ldots + x_N^2$. 

\medskip

Montrer que $f$ admet au moins un minimiseur global dans $\eR^N$.
\end{exercice}

\medskip

\begin{exercice}

Soit $f : \eR^N \rightarrow \eR$ une fonction \emph{strictement convexe}, autrement dit, pour tout couple $(x,x')$ de vecteurs de $\eR^N$, avec $x\neq x'$, et pour tout $\theta \in ]0;1[$, on vérifie :
\[
f( \theta x + (1-\theta) x' ) < \theta f(x) + (1-\theta) f(x').
\]

\begin{enumerate}

\item Une fonction strictement convexe est-elle convexe ? Réciproquement, une fonction convexe est-elle toujours strictement convexe ?

\item Montrer qu'une fonction strictement convexe admet au plus un seul minimiseur global.

\end{enumerate}

\end{exercice}

\newpage

\section{Fonctions à plusieurs variables (II)}

On va s'intéresser ici à la minimisation de fonctions, et voir comment approcher le minimiseur d'une fonction.

\medskip

Un exemple pratique de problème de minimisation est le suivant : on se donne $m$ points du plan de coordonnées $(x_i,y_i)$, $i=1,\ldots,m$ et on cherche la droite qui approche le nuage de points \emph{au sens des moindres carrés}. Autrement dit on cherche les coefficients $\alpha$ et $\beta$ qui minimisent la fonction :
\begin{equation} \label{moindre_carres}
J(\alpha,\beta) = \sum_{i=1}^m ( y_i - \alpha - \beta x_i )^2
\end{equation}
appelée \emph{erreur quadratique minimale}.

\medskip

De façon générale, pour $N \geq 1$, on notera 
\[
J : \eR^N \rightarrow \eR
\]
la fonction à minimiser, et on notera aussi son gradient :
\[
\nabla J = \left [ \frac{\partial J}{\partial x_1} , 
\frac{\partial J}{\partial x_2} , \ldots , \frac{\partial J}{\partial x_N} \right ]^\mathrm{T}.
\]
La notation $\| x \|_2^2 = x_1^2 + \ldots + x_N^2$ désignera toujours la norme euclidienne dans $\eR^N$.

\medskip

On dira que $d \in \eR^N$ est une \emph{direction de descente} de $J$ en $x\in \eR^N$ s'il existe $\delta > 0$ tel que pour tout $t \in [0,\delta]$, on vérifie
\[
J ( x + t d ) \leq J(x).
\]
La direction est dite \emph{stricte} si l'inégalité ci-dessus est stricte pour $t>0$.

\medskip

\begin{exercice}

On supposera que $J$ est différentiable au point $x$, et que $d$ est une direction de descente de $J$ en $x$.

\begin{enumerate}

\item Montrer que \[ ( \nabla J (x) , d ) \leq 0 \]
où $(\cdot,\cdot)$ est le produit scalaire dans $\eR^N$.

\item Montrer que si, de plus, $\nabla J (x) \neq 0$, alors $d = -\nabla J(x)$ est une direction de descente stricte de $J$ en $x$.

\end{enumerate}

\end{exercice}

\medskip

Soit $\lambda > 0$ un paramètre, qu'on appelle \emph{le pas}. On appelle \emph{algorithme de gradient à pas fixe} la méthode itérative suivante :
\begin{equation}\label{eq:gradient}
x^{k+1} = x^k - \lambda \nabla J (x^k),
\end{equation}
en partant d'un point $x^0 \in \eR^N$. Cette méthode itérative permet d'obtenir une suite $(x^k)$ susceptible d'approcher, sous conditions, un minimiseur de la fonction $J$.

\medskip

\begin{exercice}

Montrer que la méthode \eqref{eq:gradient} peut s'interpréter comme une méthode d'approximation de point fixe, pour une fonctionnelle qu'on explicitera. Quand $N=1$, quelle méthode déjà connue retrouve-t'on ?

\end{exercice}

\medskip

\begin{exercice}

Appliquer la méthode de gradient à pas fixe pour retrouver le minimum de 
\[
J(x_1,x_2)=x_1^2+3x_2^2+2x_1x_2-2x_1-10x_2+9.
\]
\end{exercice}

\medskip

\begin{exercice}

Le but de cet exercice est d'établir un résultat de convergence pour la méthode du gradient à pas fixe. On va supposer pour cela que $J$ est différentiable, et aussi qu'elle est \emph{fortement convexe} de paramètre $\alpha>0$, autrement dit que pour tout $x,y \in \eR^N$: %, pour tout $\theta \in [0,1]$:
\[
%J ( \theta x + (1-\theta) y ) \leq \theta J(x) + ( 1 - \theta) J(y) - \alpha %\frac{\theta (1-\theta)}{2} \| x - y \|_2^2.
( \nabla J(x) - \nabla J(y) , x - y ) \geq \alpha \| x - y \|^2_2.
\]
On suppose aussi que l'application $\nabla J : \eR^N \rightarrow \eR^N$ est lipschitzienne, autrement dit qu'il existe $L>0$ tel que pour tout $x,y \in \eR^N$:
\[
\| \nabla J (x) - \nabla J (y) \|_2 \leq L \| x - y \|_2.
\]
On va supposer finalement : 
\begin{equation}\label{eq:valeur_alpha}
0 < \lambda < \frac{2\alpha}{L^2}.
\end{equation}

\begin{enumerate}

\item Montrer d'abord que $J$ admet un unique minimiseur, noté $x^*$.

\item Soit $J_\lambda : \eR^N \rightarrow \eR^N$ définie par 
\[
J_\lambda ( x ) = x - \lambda \nabla J (x),
\]
pour tout $x\in\eR^N$.
Montrer que, pour tout $x,y\in\eR^N$  :
\[
\| J_\lambda (x) - J_\lambda(y) \|^2_2 \leq (1 - 2 \lambda \alpha + \lambda^2 L^2)  \| x - y \|^2_2.
\]

\item En utilisant \eqref{eq:valeur_alpha}, en déduire qu'il existe $\rho \in ] 0 ; 1 [$ tel que 
\[
\| J_\lambda (x) - J_\lambda(y) \|_2 \leq \rho \| x - y \|_2.
\]

\item Montrer alors que la suite $(x^k)$ donnée par \eqref{eq:gradient} vérifie
 pour tout $k\geq0$ :
 \[
  \| x^* - x^{k+1} \|_2 \leq \rho \| x^* - x^k \|_2.
 \]
 
\item En déduire que la suite $(x^k)$ tend vers $x^*$.
 
\end{enumerate}

\end{exercice}

\medskip

\begin{exercice}

Soit une matrice $A \in \eR^{N,M}$ avec $M<N$. 

\medskip

Pour $y\in\eR^N$, on introduit la fonction $J : \eR^M \rightarrow \eR$ définie par
\begin{equation} \label{eq:moindres_carres2}
J(x) = \| A x - y \|^2_2,
\end{equation}
où $x \in \eR^M$.

\begin{enumerate}

\item Que peut-on dire sur le rang de $A$ ? En particulier, le système 
\[
A x = y
\]
admet-il toujours une solution pour $y\in\eR^N$ ?

\item Montrer que la fonctionnelle $J$ admet un unique minimiseur.

\item Montrer que le minimiseur de $J$ est solution d'un système linéaire que l'on explicitera.

\item Peut-on retrouver la fonctionnelle \eqref{moindre_carres} comme un cas particulier de \eqref{eq:moindres_carres2} ? Si oui, expliciter la matrice $A$ correspondante.

\end{enumerate}

\end{exercice}

%\begin{exercice}
%
%Dérivées et calculs d'extrema : exercices ISIFC et L1 Eco.
%
%\end{exercice}

%\medskip

%\begin{exercice}
%
%Le théorème de Schwartz : exemples et contre-exemples ?
%
%\end{exercice}

%\medskip

%\begin{exercice}

%Un peu de descente de Gradient. Voir Ern et Stoltz, et peut-être Allaire aussi.

%\end{exercice}

% VOIR AUSSI EXOS PREPA.

\newpage

\section{Courbes paramétrées}

Dans ce dernier TD, nous allons étudier des notions sur les courbes paramétrées. Soit $D \subset \eR$, une \emph{courbe paramétrée plane} est une application
\[
f : \left \{
\begin{array}{lcl}
D & \rightarrow & \eR^2 \\
t & \mapsto & f(t) = (x(t),y(t)).
\end{array}
\right .
\]

\medskip

\begin{exercice}

Une \emph{courbe de Lissajous} est donnée par les équations:
\[
\left \{
\begin{array}{lcl}
x(t) & = & \sin(2t) \\
y(t) & = & \sin(3t),
\end{array}
\right .
\]
pour $t\in\eR$.

\begin{enumerate}
\item Déterminer un domaine d'étude le plus simple possible pour cette courbe.

\item Dans quelle partie du plan se situe la courbe ?

\item Trouver les points où la tangente à la courbe est verticale, puis horizontale.

\item Tracer la courbe.

\end{enumerate}

\end{exercice}

\medskip

\begin{exercice}

Trouver les points multiples de la courbe :
\[
\left \{
\begin{array}{lcl}
x(t) & = & 2t + t^2 \\
y(t) & = & 2t - \frac1{t^2},
\end{array}
\right .
\]
pour $t\in\eR^*$.

\end{exercice}

\medskip

\begin{exercice}

Soit $(\cdot,\cdot)$ la notation pour le produit scalaire dans $\eR^2$, et $\| \cdot \|_2$ celle pour la norme associée.
Soit $D \subset \eR$ et deux applications $f$ et $g$ définies sur $D$ et à valeurs dans $\eR^2$. Soit $t_0 \in D$, tel que $f$ et $g$ soient dérivables en $t_0$.

\begin{itemize}

\item Montrer que l'application 
\[
\phi(t) = ( f(t) , g(t) )
\]
est dérivable en $t_0$ et que 
\[
\phi'(t_0) = \left ( \frac{d}{dt} f(t_0) , g(t_0) \right ) + \left ( f(t_0) , \frac{d}{dt}g(t_0) \right ).
\]

\item En supposant que $f(t_0) \neq 0$, montrer que l'application
\[
t \rightarrow  \| f(t) \|_2
\]
est dérivable, et que sa dérivée est
\[
\frac{d}{dt} \| f(t) \|_2 (t_0) = \frac{1}{\| f(t_0) \|_2} 
\left ( f(t_0) , \frac{d}{dt} f(t_0) \right ) .
\]

\end{itemize}

\end{exercice}

\medskip

\begin{exercice}

Soit $t \rightarrow f(t)=( \cos(t) , \sin(t) )$ une paramétrisation du cercle de centre $O$ et de rayon $1$. Montrer en utilisant l'exercice précédent qu'en tout point $f(t)$ du cercle, la tangente est orthogonale au rayon.

\end{exercice}

\medskip

\begin{exercice}

Etudier le point singulier à l'origine de :
\[
\left \{
\begin{array}{lcl}
x(t) & = & t^2 \ln(1+t) \\
y(t) & = & t^2 (e^{t^2} - 1)
\end{array}
\right . .
\]

\end{exercice}

\medskip

\begin{exercice}

Construire le \emph{folium de Descartes} d'équation
\[
x^3 + y^3 - 3 a x y = 0
\]
où le paramètre $a$ est un réel strictement positif.

\end{exercice}

\end{document}

